% TOP_PROBLEM: {"id": 5500015, "title": "Output-sensitive Convex Hull in $\\mathbb{R}^d$", "category_id": 6, "category": {"id": 6, "name": "geometry", "display_name": "Geometry", "description": "Euclidean and non-Euclidean geometry, geometric structures.", "slug": "geometry", "order_index": 6, "created_at": "2026-07-31T15:26:25.671Z"}, "status": "partially_solved", "problem_number": "AMR-054-0015", "set_id": 13}
% FIRST_POSTED: 2026-10-01
% ENTRY_KIND: statement_only
% UnsolvedMath ID 5500015; source catalogue code AMR-054-0015
% !TeX program = lualatex
% Definition and sources only; this file is not an LLM solution attempt.
\documentclass[11pt,a4paper]{article}
\usepackage[margin=25mm]{geometry}
\usepackage{fontspec}
\setmainfont{lmroman10-regular.otf}[BoldFont=lmroman10-bold.otf,
 ItalicFont=lmroman10-italic.otf,BoldItalicFont=lmroman10-bolditalic.otf]
\usepackage{amsmath,amssymb,mathtools}
\usepackage{unicode-math}
\setmathfont{latinmodern-math.otf}
\AtBeginDocument{\let\setminus\smallsetminus}
\usepackage{xurl,hyperref}
\hypersetup{colorlinks=true,urlcolor=blue}
\setcounter{secnumdepth}{0}
\setlength{\parindent}{0pt}
\setlength{\parskip}{5pt}
\setlength{\emergencystretch}{3em}
\begin{document}
\section*{Output-sensitive Convex Hull in $\mathbb{R}^d$}\label{5500015}
{\small UnsolvedMath ID 5500015 · AMR-054-0015 · Geometry · AMR Open Problem Lists}
\subsection{Definitions and mathematical statement}
Fix a dimension $d\ge4$. Given $n$ points $S\subset\mathbb R^d$, compute
the boundary facets of $P=\operatorname{conv}(S)$. In the general-position
case, $P$ is full-dimensional and every facet has $d$ vertices, so the
output consists of $f$ constant-size facet descriptions. The dimension
is fixed independently of $n$.

An output-sensitive algorithm has a running time bounded in terms of
both $n$ and the actual number $f$ of facets, rather than the largest
possible hull size for $n$ points. Determine the optimal such bound.
The natural sharp target in the geometric comparison model is
\[
                         O(n\log(f+1)+f).
\]
An algorithm should produce the facets and their supporting data, and
its accounting must include all work spent on intermediate faces or
intermediate hulls. A small final output does not justify ignoring a
large intermediate enumeration.

For degenerate inputs, a non-simplicial facet can have many vertices;
then output size should include the chosen incidence representation.
The displayed formulation isolates the general-position core of the
source question. Arithmetic operations and orientation tests on fixed-size
coordinate tuples have unit cost; coordinate bit complexity is a
separate model choice.
\subsection{Short English statement}
Can a high-dimensional convex hull be computed in time governed by
the input and its actual facet output, matching the natural
$n\log f+f$ target?
\subsection{Sources}
\begin{itemize}
\item[S1] ULAM.AI and the UnsolvedMath Contributors, supplied problems.json, record 5500015 (AMR-054-0015), AMR Open Problem Lists. Catalogue identity and source transcription; CC BY 4.0.
\url{https://huggingface.co/datasets/ulamai/UnsolvedMath}.
\item[S2] Open Problems Project - Computational Geometry, Problem 15. Original source indexed by AMR-054-0015.
\url{https://topp.openproblem.net/p15}.
\end{itemize}
\end{document}
